\documentclass{article}
\usepackage{mathtools} 
\usepackage{fontspec}
\usepackage[UTF8]{ctex}
\usepackage{amsthm}
\usepackage{mdframed}
\usepackage{xcolor}
\usepackage{amssymb}
\usepackage{amsmath}


% 定义新的带灰色背景的说明环境 zremark
\newmdtheoremenv[
  backgroundcolor=gray!10,
  % 边框与背景一致，边框线会消失
  linecolor=gray!10
]{zremark}{说明}

% 通用矩阵命令: \flexmatrix{矩阵名}{元素符号}{行数}{列数}
\newcommand{\flexmatrix}[4]{
  \[
  #1 = \begin{pmatrix}
    #2_{11}     & #2_{12}     & \cdots & #2_{1#4}   \\
    #2_{21}     & #2_{22}     & \cdots & #2_{2#4}   \\
    \vdots      & \vdots      & \ddots & \vdots     \\
    #2_{#31}    & #2_{#32}    & \cdots & #2_{#3#4}
  \end{pmatrix}
  \]
}

% 简化版命令（默认矩阵名为A，元素符号为a）: \quickmatrix{行数}{列数}
\newcommand{\quickmatrix}[2]{\flexmatrix{A}{a}{#1}{#2}}

\begin{document}
\title{注释 6.1}
\author{张志聪}
\maketitle

\begin{zremark}
  推论3（导数极限定理）的单侧情况：

  设函数$f$在点$x_0$的某右领域$U_+(x_0)$上连续，在$U_+^{\circ}(x_0)$上可导，
  且极限$\lim\limits_{x \to x_0^+} f'(x)$存在，则$f$在点$x_0$处右导数存在，且
  \begin{align*}
    f'_+(x_0) = \lim\limits_{x \to x_0^+} f'(x)
  \end{align*}
\end{zremark}

\textbf{证明：}

由书中推论3的证明过程可知，命题的条件要保证拉格朗日定理的条件成立。
因为书中的证明就是按照单侧的情况来证明的，这里无需赘述了。

\begin{zremark}
  导函数没有第一类间断点：\\
  设$f(x)$在区间$I$上可导， 证：$f'(x)$在$I$上没有第一类间断点。
\end{zremark}

\textbf{证明：}

反证法，假设$f'(x)$在$I$上有第一类间断点，即存在$x_0 \in I$，有
\begin{align*}
  f'(x_0 - 0) \neq f'(x_0 + 0)
\end{align*}
由推论3（导数极限定理）（准确的说，是以单侧的角度来看）可知
\begin{align*}
  f'_+(x_0) = f'(x_0 + 0) \\
  f'_-(x_0) = f'(x_0 - 0) 
\end{align*}
已知$f$在$I$上可导，则$f$在$x_0$处可导，所以
\begin{align*}
  f'(x_0) = f'_+(x_0) = f'_-(x_0)
\end{align*}
于是，我们有
\begin{align*}
  f'(x_0 + 0) = f'(x_0 - 0) 
\end{align*}
即$f'(x)$在$x_0$处连续，存在矛盾。

\begin{zremark}
  定理6.4 的证明
\end{zremark}

\textbf{证明：}

不妨设$f$严格单增。

\begin{itemize}
  \item 必要性

        (i)由定理6.3可知，命题成立。

        (ii) 由定理6.2推论1可知，命题成立。

  \item 充分性

        由(i)可得，$f$是递增的。
        对任意$x_1, x_2 \in (a, b)$，如果$f(x_1) = f(x_2)$，
        由于$f$是递增的，所以在区间$[x_1, x_2]$上，$f$只能是常数，
        于是在区间$[x_1, x_2]$上，$f'(x) \equiv 0$，与(ii)矛盾。

        综上，$f$是严格单增函数。
\end{itemize}

\begin{zremark}
  定理6.4推论下的“注”
  \begin{itemize}
    \item
          (1) 函数$f$在区间$(a, b)$上严格单增，如果$f$在$a$点右连续，那么$f$在$[a, b)$上严格单增。
  \end{itemize}
\end{zremark}

\textbf{证明：}

\begin{itemize}
  \item (1)

        反证法，
        (i)假设存在$x_0 \in (a, b)$，使得$f(x_0) = f(a)$，
        于是在区间$(a, x_0]$上，$f'(x) \equiv 0$，与定理6.4矛盾。

        (ii)假设存在$x_0 \in (a, b)$，使得$f(x_0) < f(a)$，
        因为$f$在$a$点右连续，由保号性可知，
        存在$x \in (a, x_0)$，使得$f(x_0) < f(x) < f(a)$，
        这与$f$严格单调矛盾。
        综上，$f$在$[a, b)$上严格单增。

\end{itemize}


\end{document}